\chapter{Caveat Emptor --- Limitations of Local Sensitivity Analysis}

\begin{flushright}
\textit{Local sensitivity is a linear truth, and nature is rarely linear for long.}
\end{flushright}

\bigskip

In 2019, Saltelli et al.~\cite{saltelli2019why} reviewed sensitivity analyses
published in 280 highly cited papers across top environmental modeling journals.
More than half used the wrong method for the question being asked.
Specifically, studies that asked global questions, how does output
uncertainty relate to input uncertainty across the full parameter range?;
used local methods: derivatives computed at a single nominal parameter value.
The derivatives were computed correctly. But they answered a different question
than the one the authors claimed to answer.

The deeper problem is not computational but epistemological: sensitivity
analysis is still often defined by method rather than purpose
\cite{razavi2021future}.
Practitioners reach for the method from their own disciplinary tradition
rather than the one best suited to their question.
A local method applied to answer a global question will produce results
that are internally correct but wrong for the stated purpose.
Before choosing an SA method, two questions must be answered: what is the
goal --- factor prioritization (which inputs most reduce output variance?),
factor fixing (which inputs can be frozen without changing the output?),
or factor mapping (which inputs drive the output above a threshold?) ---
and does the method's assumptions match the model's behavior?

\medskip
\begin{center}
\textit{Assuming local sensitivity holds globally is the triumph\\
of mathematical hope over nonlinear reality.}
\end{center}
\medskip

This chapter asks the question that should be asked before any SA is begun:
is local SA the right tool for this question?

%%------------------------------------------------------------
\section{Method-Question Mismatch}\index{method-question mismatch}
\label{sec:mismatch}
%%------------------------------------------------------------

Local sensitivity analysis is built on a derivative. A derivative
is a local object: it measures the behavior of a function at one point,
and it is only meaningful in the vicinity of that point.
This makes local SA powerful for some questions and completely wrong
for others.

Local SA can answer: \emph{which parameter most affects the model
near the nominal parameter values?}

Local SA cannot answer: \emph{which parameter most affects the model
across its full range of uncertainty?} That question requires the parameter
to actually vary across its range, and a derivative does not do that.

Local SA cannot answer: \emph{how does the model behave at parameter
values far from the nominal?} The derivative approximates linear behavior;
real models are often highly nonlinear.

The principle is simple and unforgiving. The method must match the question.
Using local SA to answer a global question is not merely suboptimal;
it can produce a ranking of parameter importance that is entirely wrong.
A parameter that appears unimportant locally may dominate the output
variance when varied across its full uncertainty range.
A parameter that appears important locally may have negligible effect
at realistic perturbation sizes if the model is highly nonlinear
in that parameter.

A small sensitivity index does not mean a parameter can be ignored globally.
It means the model is approximately linear in that parameter near the nominal
value, and the linear coefficient happens to be small.
Section~\ref{sec:globalneeded} shows a specific example where the local
ranking is the opposite of the global ranking.

\begin{selfcheck}
A published SA study reports: ``All sensitivity indices are less than 0.05,
confirming that the model is robust to parameter uncertainty.''
The study used nominal parameter values from the literature, with
no information on parameter uncertainty ranges.
Identify two specific claims in this statement that go beyond what
local SA can actually establish.
\end{selfcheck}
\begin{scAnswer}
(1) ``Robust to parameter uncertainty'': local SA says nothing about
parameter uncertainty. It describes behavior near the nominal value.
If the uncertainty range for some parameters is large (e.g., $\pm50\%$),
the model could be highly sensitive to that parameter globally even
though the local SI is small.

(2) The conclusion depends implicitly on the nominal parameter values
chosen. If the nominal values are changed, the local SI values change
too (for nonlinear models). A study with different nominal values
might find very different rankings. This sensitivity of local SA
to the choice of nominal values is itself a limitation that the study
does not address.
\end{scAnswer}

%%------------------------------------------------------------
\section{Bifurcations and SA Failure}\index{bifurcation!SA failure near}
\label{sec:bifurcations}
%%------------------------------------------------------------

Local SA requires the model output to be differentiable with respect
to the parameters at the nominal values. When this fails, local SA
fails completely. The most important way it fails is at a bifurcation.

A bifurcation is a parameter value at which the qualitative behavior
of the model changes: two equilibria merge into one, a stable fixed
point becomes unstable, or an epidemic threshold is crossed~\cite{strogatz2018nonlinear,kuznetsov2004elements}.
At a bifurcation point, the derivative $\partial u/\partial p$ is
either undefined or infinite. Local SA breaks down entirely.
The precise mathematical reason is given by the Fredholm alternative
(Appendix~\ref{appsec:solvability}): near a bifurcation the Jacobian
becomes singular, so the FSE $D_{\vec{u}}[\vec{f}]\,(\partial\vec{u}/\partial p) = \vec{r}$
may have no solution or infinitely many.
For models with periodic limit cycles, sensitivity near a bifurcation
of the periodic orbit involves Floquet multipliers; this is treated
in Appendix~\ref{app:floquet}.

\subsection{The Quadratic Example}

Consider the scalar equation $f(u; p) = u^2 + p = 0$.
For $p < 0$, there are two real roots: $u = \pm\sqrt{-p}$.
At $p = 0$, the two roots merge. For $p > 0$, there are no real roots.

The FSE for the sensitivity of the positive root $u^* = \sqrt{-p}$ is:
\[
  \pd{u^*}{p} = \frac{-\pd{f}{p}}{\pd{f}{u}} = \frac{-1}{2u^*}
              = \frac{-1}{2\sqrt{-p}}.
\]
The normalized SI is:
\[
  \SI{p}{u^*} = \frac{p}{u^*}\,\pd{u^*}{p}
              = \frac{p}{\sqrt{-p}}\cdot\frac{-1}{2\sqrt{-p}}
              = \frac{1}{2}.
\]
This is constant (independent of $p$) for all $p < 0$;
the positive root is always in exact proportion $1/2$ to any
fractional change in $p$.

But as $p \to 0^-$:
\[
  \pd{u^*}{p} = \frac{-1}{2\sqrt{-p}} \to -\infty.
\]
The derivative diverges at the bifurcation. Near $p = 0$, the
model is infinitely sensitive to the parameter $p$: any small
positive perturbation to $p$ destroys the equilibrium entirely.

\unifyingidea{3} This is Unifying Idea~3 in its most dramatic form.
The sensitivity index is the derivative-based linearization of the
model response. At a bifurcation, the function $u^*(p)$ has
infinite curvature, its graph has a vertical tangent, and
no linear approximation is valid.
The linearization does not just become inaccurate at the bifurcation;
it becomes undefined.

\unifyingidea{1} The failure at a bifurcation is a failure of the
normalized derivative. The sensitivity index $\SI{p}{q}$ exists only
when $\partial q/\partial p$ is finite. Near a bifurcation it is not,
and the index either blows up or vanishes depending on which side
of the threshold you are computing from.

\unifyingidea{2} The adjoint SA system inherits exactly the same
failure mode. The adjoint ODE is driven by the Jacobian $J_F^T$
of the forward model. Near a bifurcation, the Jacobian has a
near-zero eigenvalue, and the adjoint solution grows without bound
over the integration interval. Both the forward and adjoint
approaches fail at the same bifurcation for the same reason:
the transpose of a singular matrix is also singular.

\begin{selfcheck}
For $f(u;p) = u^2 + p$ with $p = -4$:
(a) Find the two roots $u^*_\pm$.
(b) Compute $\SI{p}{u^*_+}$ analytically. Compare with the derivative $\partial u^*_+/\partial p$.
(c) Estimate $\partial u^*_+/\partial p$ numerically using a centered difference
    with $h = 0.1, 0.01, 0.001$. At what value of $p$ near zero does
    the finite-difference estimate become unreliable?
\end{selfcheck}
\begin{scAnswer}
(a) $u^*_\pm = \pm\sqrt{4} = \pm 2$.
(b) $\partial u^*_+/\partial p = -1/(2\sqrt{-p}) = -1/4$.
$\SI{p}{u^*_+} = (p/u^*_+)(-1/(2\sqrt{-p})) = (-4/2)(-1/4) = 1/2$.

(c) Numerically at $p = -4$: all three step sizes give $\approx -0.250$ \checkmark.
As $p \to 0^-$, the function $u^*_+ = \sqrt{-p}$ has a square-root singularity;
finite differences become inaccurate when the curvature term dominates the
step-size error, which occurs when $h \gtrsim |p|$. Near $p = 0$, even very
small $h$ is not small compared to $|p|$, so finite differences fail.
\end{scAnswer}

\subsection{The SIR Model at the Epidemic Threshold}

The SIR model has a transcritical bifurcation at $R_0 = 1$:
for $R_0 < 1$, the only equilibrium is the disease-free state;
for $R_0 > 1$, the endemic equilibrium appears.

The endemic equilibrium fraction
$i^* = 1 - 1/R_0$
exists only for $R_0 > 1$. Its sensitivity with respect to $R_0$ is:
\[
  \pd{i^*}{R_0} = \frac{1}{R_0^2}, \qquad
  \SI{R_0}{i^*} = \frac{R_0}{i^*}\cdot\frac{1}{R_0^2}
                = \frac{1}{R_0(1-1/R_0)} = \frac{1}{R_0 - 1}.
\]
As $R_0 \to 1^+$, $\SI{R_0}{i^*} \to +\infty$: near the epidemic
threshold, the endemic equilibrium is infinitely sensitive to changes
in $R_0$. This is exactly the bifurcation behavior identified above.

SA of the endemic equilibrium for $R_0$ values close to 1 must be
interpreted with extreme caution. A small change in any parameter
that affects $R_0$ could eliminate the endemic equilibrium entirely
or cause it to surge. Local SA correctly identifies this extreme
sensitivity but cannot predict the qualitative change.

For $R_0 \leq 1$: SA of the endemic equilibrium is undefined;
the equilibrium does not exist. Computing SA under the assumption
that $i^*$ exists when $R_0 < 1$ is a mathematical error.
MATLAB will simply return a negative or complex value without warning.

\medskip\noindent
\textbf{Operational rule: when to stop reporting SI as a ranking.}
A sensitivity index is meaningful as a ranking tool only when
the linearization is a fair local approximation.
Two practical checks:

\begin{enumerate}[nosep]
\item \textit{Magnitude check.} If $|\SI{p}{q}| > 10$ for a normalized
      index, treat the result as a qualitative warning
      (``this parameter region is highly sensitive'') rather than a
      quantitative ranking. An SI of 100 does not mean the parameter
      matters 10$\times$ more than one with SI 10;
      it may mean the model is near a threshold.
\item \textit{Condition-number check.} Compute $\kappa(J_F)$
      (the condition number of the model Jacobian at the operating point).
      If $\kappa(J_F) > 10^4$, the sensitivity computation is
      ill-conditioned: small errors in the model evaluation
      may be amplified by $\kappa(J_F)$ in the sensitivity values.
      Always report $\kappa(J_F)$ alongside SI values computed
      near a bifurcation or near parameter degeneracy.
\end{enumerate}

%%------------------------------------------------------------
\section{Ill-Conditioning}\index{ill-conditioning}
\label{sec:illconditioning}
%%------------------------------------------------------------

For linear algebraic models, SA failure takes a different form.
When the coefficient matrix $A$ in $A\vec{u} = \vec{b}$ is
nearly singular, small changes in the right-hand side $\vec{b}$
(or in the matrix entries) produce large changes in the solution.
The sensitivity indices are large not because the model is interesting
but because the model is poorly identified.

\subsection{The Condition Number}

The \textbf{condition number}\index{condition number} $\kappa(A) = \|A\|\|A^{-1}\|$
(see Appendix~\ref{appsec:SVD} for the SVD-based definition $\kappa_2 = \sigma_1/\sigma_r$)
measures how much the solution $\vec{u}$ can change relative to
a change in the data. For the 2-norm,
\begin{equation}
  \kappa_2(A) = \frac{\sigma_{\max}(A)}{\sigma_{\min}(A)},
  \label{eq:condnum}
\end{equation}
where $\sigma_{\max}$ and $\sigma_{\min}$ are the largest and smallest
singular values of $A$. A matrix with $\kappa_2(A) = 1$ is perfectly
conditioned (orthogonal); a matrix approaching singularity has
$\sigma_{\min} \to 0$ and $\kappa_2(A) \to \infty$.

The key bound: if the data changes by a relative amount $\epsilon$,
the solution changes by at most $\kappa(A)\,\epsilon$:
\[
  \frac{\|\delta\vec{u}\|}{\|\vec{u}\|} \leq \kappa(A)\,
  \frac{\|\delta\vec{b}\|}{\|\vec{b}\|}.
\]
When $\kappa(A) \gg 1$, even tiny changes in the data produce
large changes in the solution.

\subsection{SA and Ill-Conditioning}

For SA, the key insight is that large sensitivity indices are not
necessarily a property of the model's \emph{true} behavior.
They may simply reflect ill-conditioning: the fact that the linear
system is nearly singular.

A colleague who reports ``the sensitivity of $u_1$ to $b_1$
is extremely large'' may be reporting either:
(a) a genuine model property (the output truly varies strongly
    with this parameter), or
(b) a numerical artifact (the linear system is nearly singular
    and small input changes have large numerical effects on the solution).

Distinguishing (a) from (b) requires computing $\kappa(A)$.
If $\kappa(A) \approx 10^6$, then numerical errors of size
$\varepsilon_{\text{mach}} \approx 10^{-16}$ can contaminate the
solution at the level of $10^{-10}$: 6 digits of the solution
may be meaningless.

Ill-conditioning is itself a SA result: it warns the investigator
that the model parameters are nearly unidentifiable from data,
and that predictions cannot be trusted without more precise
parameter measurements.
This is not a failure of SA; it is one of the most important
things SA can reveal.

\begin{selfcheck}
The $2\times2$ matrix $A(\epsilon) = \begin{pmatrix} 1 & 1 \\ 1 & 1+\epsilon \end{pmatrix}$
with $\epsilon > 0$:
(a) Compute $\det(A(\epsilon))$ and show $A \to$ singular as $\epsilon \to 0$.
(b) Compute the condition number $\kappa_2(A)$ for $\epsilon = 1, 0.1, 0.01$.
(c) For $\vec{b} = (2, 2+\epsilon)^T$, solve $A\vec{u} = \vec{b}$ and compute
    $\SI{\epsilon}{u_1}$. What happens to this SI as $\epsilon \to 0$?
\end{selfcheck}
\begin{scAnswer}
(a) $\det = \epsilon \to 0$ as $\epsilon \to 0$. Singular at $\epsilon = 0$.
(b) $\kappa_2(A) = \sigma_{\max}/\sigma_{\min}$. For $\epsilon = 1$: $\kappa \approx 5.83$.
For $\epsilon = 0.1$: $\kappa \approx 40$. For $\epsilon = 0.01$: $\kappa \approx 400$.
As $\epsilon \to 0$: $\kappa \to \infty$.

(c) $A^{-1} = \epsilon^{-1}\begin{pmatrix}1+\epsilon & -1\\ -1 & 1\end{pmatrix}$.
$\vec{u} = A^{-1}\vec{b} = (1, 1)^T$ for all $\epsilon > 0$.
$\SI{\epsilon}{u_1}$: differentiate $u_1 = (2(1+\epsilon) - (2+\epsilon))/\epsilon
= \epsilon/\epsilon = 1$ with respect to $\epsilon$, but the simplification
$u_1 = 1$ holds exactly, so $\SI{\epsilon}{u_1} = 0$. The sensitivity vanishes!

This illustrates a subtlety: large condition number does not always imply large SI.
What the condition number warns is that \emph{numerical} computation of $\vec{u}$
will be inaccurate; the exact mathematical SI may nonetheless be finite or zero.
\end{scAnswer}

%%------------------------------------------------------------
\section{Separatrices and Trajectory Sensitivity}\index{separatrix}\index{trajectory sensitivity}
\label{sec:separatrices}
%%------------------------------------------------------------

Dynamical systems often have regions of qualitatively different behavior
separated by a \textbf{separatrix}: a boundary in state space that
trajectories cannot cross. SA computed on one side of a separatrix
describes behavior qualitatively different from SA on the other side.

For the SIR model, the relevant separatrix is the epidemic threshold
$S(0) > \gamma N / (k\beta)$ (equivalently, $R_0 > 1$). If the
initial susceptible population exceeds this threshold, an epidemic
occurs. If it does not, the infection dies out without spreading.

Suppose SA is performed at a parameter value where an epidemic occurs
($R_0 = 1.5$) and the sensitivity index $\SI{\beta}{I_{\max}}$
is computed. This index describes how the epidemic peak changes with
a small perturbation to $\beta$ within the epidemic regime.
It says nothing about what happens if $\beta$ is reduced below the
threshold: in that regime, $I_{\max}$ is not a smooth function of $\beta$
it drops to zero discontinuously as the separatrix is crossed.

The practical implication: if any parameter perturbation of interest
crosses a qualitative threshold (endemic vs.\ disease-free, outbreak
vs.\ no outbreak, stable vs.\ oscillatory), local SA cannot be used to
predict the response. The model must be analyzed separately in each
qualitative regime, and the boundary itself requires special treatment.

A related limitation applies to models with multiple stable attractors.
When a model can settle into qualitatively different long-run states
depending on initial conditions, the sensitivity index computed near one
attractor says nothing about behavior near another.
Extensions of the SIR model that exhibit bistability (where both a
disease-free and an endemic equilibrium are stable for the same parameter
values) require separate SA studies for each attractor.
Local SA at one equilibrium cannot warn the analyst that another
equilibrium exists.

%%------------------------------------------------------------
\section{When Global SA Is Needed}\index{global sensitivity analysis!when needed}
\label{sec:globalneeded}
%%------------------------------------------------------------

\begin{flushright}
\textit{\small Local SA studies the neighborhood. Global SA studies the country.}
\end{flushright}

Local SA fails to answer the right question in three specific
situations. When any of these conditions holds, Chapter~9's global
SA methods are required. The practical decision framework for
choosing between local and global SA is developed at length by
Saltelli et al.~\cite{saltelli2004sensitivity,saltelli2006sensitivity}.

To reiterate the misconception that opened this chapter:
a sensitivity index near zero does not mean a parameter is globally
unimportant. It means the linearization at the nominal value is
nearly flat. What happens away from the nominal, across the
full uncertainty range of the parameter, requires global SA to determine.

\subsection{Large Parameter Uncertainty}

When a parameter is known only to within a factor of 2 or 3;
common in epidemiology, ecology, and environmental modeling;
the linearization underlying local SA may be very poor across
the plausible range.

As a concrete example, consider a model where the \QOI depends
on a parameter $p$ through $q(p) = e^p$. At nominal value $p^* = 0$:
$\SI{p}{q} = p^*(dq/dp)/q = 0 \cdot 1/1 = 0$. The local SI is zero.
But if the uncertainty range is $[-2, 2]$, the \QOI varies from
$e^{-2} \approx 0.14$ to $e^2 \approx 7.4$: a factor of 53.
The parameter is globally important but invisible to local SA.

\subsection{Highly Nonlinear Models}

Two models with identical local SIs for a parameter $p$ may
behave very differently when $p$ is perturbed by $20\%$:
one model may be nearly linear (the SI is a good approximation),
while the other may be highly curved (the SI understates the true change).

The sigma-normalized sensitivity index (Section~9.1) provides
the minimal bridge: it weights the local SI by the parameter's
standard deviation, giving a measure that at least accounts for
the relative magnitudes of parameter uncertainties.

\subsection{Correlated Parameters}

Local SA under the independence assumption gives misleading results
when parameters are correlated (Chapter~3). Global SA methods in
Chapter~9 include specific tools for handling this case.

\subsection{The Portfolio Example}

Consider a model $Y = s_1 Q_1 + s_2 Q_2$ with $s_1 = 2$, $s_2 = 1$
and parameter standard deviations $\sigma_1 = 1$, $\sigma_2 = 3$.

Local SA gives $\SI{Q_1}{Y} = 2 > \SI{Q_2}{Y} = 1$:
parameter~1 appears more important.

The sigma-normalized index:
$S_i^\sigma = (s_i\sigma_i)/\sigma_Y$,
where $\sigma_Y = \sqrt{4+9} = \sqrt{13}$.
So $S_1^\sigma = 2/\sqrt{13} \approx 0.55$ and $S_2^\sigma = 3/\sqrt{13} \approx 0.83$.

Parameter~2 dominates the output variance, even though it has a smaller
local SI. The local ranking is inverted because parameter~2 has much
larger uncertainty. Chapter~9's Sobol indices provide the rigorous
global version of this sigma-normalized assessment.

\begin{selfcheck}
For the SIR model, suppose $\beta$ has uncertainty $\sigma_\beta = 0.02$
and $\tau$ has uncertainty $\sigma_\tau = 3$~days, at nominal values
$\beta^* = 0.06$ and $\tau^* = 7$~days.

(a) From Chapter~2, $\SI{\beta}{R_0} = \SI{\tau}{R_0} = 1$.
   Which parameter appears more important from local SA alone?

(b) Compute $\sigma_\beta/\beta^*$ and $\sigma_\tau/\tau^*$ (the relative uncertainties).
   Which parameter has larger relative uncertainty?

(c) Compute the sigma-normalized indices for $R_0$ with respect to
    $\beta$ and $\tau$. Which parameter drives more variance in $R_0$?
\end{selfcheck}
\begin{scAnswer}
(a) Both have SI = 1, so local SA alone says they are equally important.

(b) $\sigma_\beta/\beta^* = 0.02/0.06 \approx 0.33$ (33\%);
$\sigma_\tau/\tau^* = 3/7 \approx 0.43$ (43\%). The mean infectious
period has larger relative uncertainty.

(c) $S_\beta^\sigma = \SI{\beta}{R_0} \cdot (\sigma_\beta/\beta^*) = 1 \times 0.33 = 0.33$.
$S_\tau^\sigma = \SI{\tau}{R_0} \cdot (\sigma_\tau/\tau^*) = 1 \times 0.43 = 0.43$.
The mean infectious period drives more variance in $R_0$, despite both
having the same local SI. Chapter~9 will quantify this more precisely
with Sobol indices.
\end{scAnswer}

%%------------------------------------------------------------
\section{When SA Guides Data Collection}
\label{sec:ch8datacollection}
%%------------------------------------------------------------

The three chapters preceding this one have identified when SA
can mislead. This section completes the practical argument
the book has been building since Chapter~1: how to use SA
\emph{correctly} to make decisions about field measurement.

\textbf{The decision chain.}
A sensitivity study produces a ranked list of parameters.
Translating that list into a field-measurement plan requires
three additional steps that SA alone cannot provide:

\begin{enumerate}
\item \textit{Identifiability:} Can the high-sensitivity parameter
      actually be estimated from the available data, or is it
      structurally confounded with another parameter?
      Section~\ref{sec:correlatedPOIs} showed that parameters
      with identical sensitivity indices are not separately identifiable
      from outbreak data alone; measuring one is equivalent
      to measuring the other.
      For the SIR model, $k$ and $\beta$ always appear as the product
      $k\beta$, so field studies should estimate $k\beta$ directly
      (e.g., from secondary attack rates) rather than $k$ and $\beta$
      separately.

\item \textit{Measurability:} Is the high-sensitivity parameter
      actually measurable with available instruments and budget?
      The mean infectious period $\tau$ appears in the top tier
      of every SIR sensitivity analysis, and it is measurable:
      clinicians record the duration from symptom onset to recovery.
      The lifespan $L$ has negligible SI for short-horizon models
      and is measurable anyway from demographic records.
      A parameter with large SI but no feasible measurement pathway
      should be \emph{reduced} rather than measured, by redesigning
      the intervention.

\item \textit{Actionability:} Can the high-sensitivity parameter
      be reduced by a feasible public health intervention?
      SA of the SIR model with $J = \max_t I(t)$ consistently identifies
      $k$ (contact rate) and $\beta$ (transmission probability) as the
      most influential parameters.
      Contact rates can be reduced by social distancing, school closures,
      and remote-work policies.
      Transmission probability per contact can be reduced by mask
      requirements, hand-hygiene campaigns, and vaccination.
      These two routes correspond to different interventions with
      different costs, compliance rates, and social acceptability.
      SA identifies both as equally valuable targets; the choice between
      them requires policy judgment beyond the model.
\end{enumerate}

\textbf{The two-budget problem.}
A common practical scenario: a government agency has funding to
(1) measure exactly $k_{\text{meas}}$ parameters precisely in the field,
and (2) implement exactly $k_{\text{int}}$ interventions.
The correct SA-guided procedure is:

\begin{enumerate}[nosep]
\item Run local SA (or Morris screening for large parameter sets)
      to rank all parameters by $|\SI{p_j}{\text{peak}~I}|$.
\item From the top-ranked parameters, apply the identifiability filter:
      remove any parameter that is structurally confounded
      with a higher-ranked parameter.
\item Apply the measurability filter: remove any parameter for which
      no feasible measurement protocol exists within the budget.
\item The remaining top $k_{\text{meas}}$ parameters are the
      measurement priorities.
\item From the ranked list, apply the actionability filter separately:
      the top $k_{\text{int}}$ actionable parameters are the
      intervention priorities. These sets may overlap.
\end{enumerate}

\textbf{What the model cannot tell you.}
The complete picture requires one more piece: the second semester.
Smith's \\emph{Uncertainty Quantification}~\\cite{smith2014uncertainty}
Chapters~7--8 covers Bayesian parameter estimation --- the step
that actually quantifies how much a given measurement budget
reduces parameter uncertainty.
SA tells you \emph{which} parameters matter most; parameter estimation
tells you \emph{how precisely} you can learn them from data;
optimal experimental design (Smith Chapter~11) tells you
\emph{how to design the measurement study} to learn those parameters
as efficiently as possible.

These three disciplines --- SA, parameter estimation, and optimal
experimental design --- form a complete pipeline for model-based
decision making. This book has equipped you for the first stage.

\medskip\noindent
\textit{If a parameter is influential and poorly known,
that is not a nuisance. That is your research agenda.}

%%------------------------------------------------------------
\section{MATLAB Laboratory: Bifurcation and Ill-Conditioning}
\label{sec:ch8matlab}
%%------------------------------------------------------------

\textbf{Part 1: The quadratic bifurcation.}
For $f(u;p) = u^2 + p$, compute the positive root $u^* = \sqrt{-p}$
and the SI $\SI{p}{u^*} = 1/2$ analytically.
Plot $u^*$ vs.\ $p$ for $p \in [-4, -0.01]$.
Estimate $\partial u^*/\partial p$ by centered differences at
$p = -4, -1, -0.1, -0.01, -0.001$.
Compare the finite-difference estimates with the analytic value.
At what value of $p$ does the finite-difference estimate become unreliable?

\textbf{Part 2: SIR bifurcation diagram using MATCONT.}
Download and install MATCONT (\url{https://matcont.sourceforge.io}).
Define the SIR equilibrium equations as a MATCONT system.
Use MATCONT's continuation capability to trace the endemic equilibrium
$i^* = 1 - 1/R_0$ as $R_0$ increases from 0.5 to 3.0.
MATCONT will automatically detect the transcritical bifurcation at $R_0 = 1$
and flag it. Plot the bifurcation diagram (equilibrium value vs.\ $R_0$).
If MATCONT is unavailable, trace the equilibrium manually:
for $R_0 \leq 1$, $i^* = 0$ (disease-free); for $R_0 > 1$, $i^* = 1 - 1/R_0$.

\textbf{Part 3: FSA comparison across $R_0 = 1$.}
Using the full SIR ODE system, compute $\SI{\beta}{I(T)}$ at
$T = 90$~days for $R_0 = 0.8, 1.0, 1.2, 2.0$.
At each value, use centered finite differences.
Tabulate the results. At $R_0 = 1$: does the SI blow up, return NaN,
or behave normally? Try $h = 0.01, 0.001, 0.0001$ to see what happens.

\textbf{Part 4: Condition number of the SIR Jacobian as a function of $R_0$.}
For the SIR endemic equilibrium, the Jacobian $J_F$ at $(S^*, I^*, R^*)$
depends on $R_0$. Compute $\kappa_2(J_F)$ using MATLAB's \texttt{cond}
for $R_0 \in \{1.01, 1.05, 1.1, 1.5, 2.0, 3.0\}$.
Plot $\kappa_2(J_F)$ vs.\ $R_0$.
Observe the spike as $R_0 \to 1^+$: the Jacobian becomes nearly singular
at the bifurcation. What does this imply about SA of the SIR endemic equilibrium
for parameter values that place $R_0$ close to 1?

\textbf{Part 5: Sigma-normalized SIs for the SIR model.}
Using the SIR model with uncertainty ranges $\sigma_k = 1$
(contacts/day), $\sigma_\beta = 0.02$, $\sigma_\tau = 3$~days,
$\sigma_L = 500$~days (all at $1\sigma$), compute the sigma-normalized
indices for $R_0$. Compare the ranking with the local SI ranking
from Chapter~2. Which parameter most affects $R_0$ given its
realistic uncertainty range?

%%------------------------------------------------------------
\section{Exercises}
%%------------------------------------------------------------

\begin{exercise}[\bloom{L3}]
\label{ex:ch8:quadratic}
For the equation $g(u; p) = u^3 - 3pu + 2 = 0$:
\begin{enumerate}[label=(\alph*)]
  \item Find all real roots at $p = 1$. Which root changes most rapidly with $p$?
  \item Compute $\SI{p}{u^*}$ for each real root at $p = 1$
        using the implicit function theorem.
  \item Plot the roots as functions of $p$ for $p \in [0.8, 1.2]$.
        At what value of $p$ do two roots merge?
  \item At the bifurcation value of $p$, what does the FSE predict?
        Verify numerically using finite differences.
\end{enumerate}
\end{exercise}

\begin{exercise}[\bloom{L3}]
\label{ex:ch8:SIRthreshold}
Using the SIR model with $k = 5$, $\tau = 7$~days, $L = 10{,}000$~days:
\begin{enumerate}[label=(\alph*)]
  \item Compute $\SI{\beta}{I(T)}$ for $\beta = 0.04, 0.06, 0.08, 0.12$
        at $T = 90$~days using centered finite differences.
        Note: $R_0 = k\beta\tau$ takes values $1.4, 2.1, 2.8, 4.2$.
  \item Now recompute for $\beta = 0.02$ ($R_0 = 0.7$).
        Compare with the results for $R_0 > 1$. Explain any qualitative differences.
  \item If you were to add a bifurcation analysis tool (e.g., MATCONT),
        what would you expect the bifurcation diagram to look like for
        the endemic equilibrium of this SIR model as $\beta$ varies?
\end{enumerate}
\end{exercise}

\begin{exercise}[\bloom{L4}]
\label{ex:ch8:conditioning}
For the $3\times3$ linear system $A\vec{u} = \vec{b}$ where:
\[
  A = \begin{pmatrix} 1 & 2 & 3 \\ 4 & 5 & 6 \\ 7 & 8 & 9+\epsilon \end{pmatrix},
  \quad \vec{b} = A\vec{u}_{\text{exact}}, \quad \vec{u}_{\text{exact}} = (1,1,1)^T.
\]
\begin{enumerate}[label=(\alph*)]
  \item For $\epsilon = 1, 0.1, 0.01, 0.001, 0.0001$, compute $\kappa(A)$
        and $\|\hat{\vec{u}} - \vec{u}_{\text{exact}}\|$ where
        $\hat{\vec{u}} = A\backslash\vec{b}$ in MATLAB.
  \item Plot $\|\hat{\vec{u}} - \vec{u}_{\text{exact}}\|$ vs.\ $\kappa(A)$
        on log-log axes. What is the approximate slope?
  \item What does this experiment tell you about interpreting large SI values
        for nearly singular models?
\end{enumerate}
\end{exercise}

\begin{exercise}[\bloom{L5}]
\label{ex:ch8:report}
You are asked to perform SA on an epidemic model whose best-fit
parameter values give $R_0 = 1.04$, very close to the epidemic threshold.
Write a one-page advisory report for a public health agency explaining:
\begin{enumerate}[label=(\alph*)]
  \item What local SA can and cannot tell you at this parameter value.
  \item What specific numerical problems to expect when computing
        sensitivity indices near $R_0 = 1$.
  \item What additional analyses you would recommend before advising
        on interventions.
  \item How you would explain the concept of a bifurcation and its
        implications to a non-mathematical policy maker.
\end{enumerate}
\end{exercise}

\begin{exercise}[\bloom{L6}]
\label{ex:ch8:critique}
A published modeling paper reports: ``Sensitivity analysis of the
endemic equilibrium of our malaria model found that all indices
are less than $0.1$, indicating the model is robust to parameter
uncertainty. The basic reproduction number at fitted parameter
values is $R_0 = 1.15$.''

Identify at least three specific methodological concerns with
this analysis and its conclusions. For each concern, state:
(a) what the problem is; (b) what evidence the authors would need
to provide to address it; and (c) what alternative analysis
would give a more complete picture.
\end{exercise}

%%------------------------------------------------------------

